\section{引言：图像分类——计算机视觉的核心任务}

本节课是 CS231N 的第二讲，在第一讲介绍计算机视觉全貌的基础上，本讲将聚焦于\textbf{图像分类}这一核心任务，并介绍两种基本的数据驱动方法：\textbf{K 近邻分类器}（K-Nearest Neighbor）和\textbf{线性分类器}（Linear Classifier）。线性分类器是构建深度神经网络的最重要基础模块。

\begin{figure}[H]
\centering
\includegraphics[width=0.95\textwidth]{slides-images/slide-005.jpg}
\caption{本讲课程大纲：图像分类任务定义，两种数据驱动方法——K 近邻和线性分类器\protect\footnotemark}
\end{figure}
\footnotetext{Slides 第6页。}

\begin{importantbox}{本讲核心目标}
\begin{enumerate}
    \item 定义图像分类任务及其挑战
    \item 理解数据驱动方法（data-driven approach）的基本范式
    \item 掌握 K 近邻分类器的原理、超参数选择与局限性
    \item 深入理解线性分类器的代数、视觉和几何三种视角
    \item 学习两种损失函数：多类 SVM 损失和 Softmax（交叉熵）损失
\end{enumerate}
\end{importantbox}

\subsection{什么是图像分类}

给定一张图像和一组\textbf{预定义标签}（如 dog, cat, truck, plane 等），图像分类的任务是为该图像指定一个正确的类别标签。

\begin{figure}[H]
\centering
\includegraphics[width=0.85\textwidth]{slides-images/slide-006.jpg}
\caption{图像分类任务定义：给定图像和候选标签集合，输出正确的类别标签\protect\footnotemark}
\end{figure}
\footnotetext{Slides 第7页。}

对人类而言，这是一个极其简单的任务——我们的认知系统天然具备对图像的整体理解能力。但对计算机来说，图像只是一个巨大的数字张量。

\subsection{语义鸿沟（Semantic Gap）}

\begin{figure}[H]
\centering
\includegraphics[width=0.95\textwidth]{slides-images/slide-007.jpg}
\caption{语义鸿沟：人类看到``猫''，计算机看到的是 $800 \times 600 \times 3$ 的整数张量，每个元素取值在 $[0, 255]$\protect\footnotemark}
\end{figure}
\footnotetext{Slides 第8页。}

一张彩色图像在计算机中表示为一个三维张量，维度为 $H \times W \times 3$（高度 $\times$ 宽度 $\times$ RGB三通道）。每个像素的每个通道是一个 0--255 之间的 8-bit 整数值。人类感知到的``猫''和计算机看到的数字矩阵之间存在巨大的\textbf{语义鸿沟}。

\begin{knowledgebox}{图像的数字表示}
图像通常使用 RGB 24-bit 格式存储：每个像素有红（R）、绿（G）、蓝（B）三个通道，每个通道 8 bit（$2^8 = 256$ 个取值），因此像素值范围为 $[0, 255]$。有时还有第四个 alpha 通道（透明度），构成 32-bit RGBA 格式。
\end{knowledgebox}

\subsection{图像分类面临的挑战}

图像分类的困难在于，即使是同一个物体，以下因素的变化都会导致像素值发生巨大变化：

\begin{figure}[H]
\centering
\begin{subfigure}[b]{0.48\textwidth}
\includegraphics[width=\textwidth]{slides-images/slide-008.jpg}
\caption{视角变化（Viewpoint Change）}
\end{subfigure}
\hfill
\begin{subfigure}[b]{0.48\textwidth}
\includegraphics[width=\textwidth]{slides-images/slide-009.jpg}
\caption{光照变化（Illumination）}
\end{subfigure}
\caption{图像分类的挑战（一）：视角和光照变化\protect\footnotemark}
\end{figure}
\footnotetext{Slides 第9--10页。}

\begin{figure}[H]
\centering
\begin{subfigure}[b]{0.48\textwidth}
\includegraphics[width=\textwidth]{slides-images/slide-010.jpg}
\caption{背景杂乱（Background Clutter）}
\end{subfigure}
\hfill
\begin{subfigure}[b]{0.48\textwidth}
\includegraphics[width=\textwidth]{slides-images/slide-011.jpg}
\caption{遮挡（Occlusion）}
\end{subfigure}
\caption{图像分类的挑战（二）：背景杂乱和遮挡\protect\footnotemark}
\end{figure}
\footnotetext{Slides 第11--12页。}

\begin{figure}[H]
\centering
\begin{subfigure}[b]{0.48\textwidth}
\includegraphics[width=\textwidth]{slides-images/slide-012.jpg}
\caption{形变（Deformation）}
\end{subfigure}
\hfill
\begin{subfigure}[b]{0.48\textwidth}
\includegraphics[width=\textwidth]{slides-images/slide-013.jpg}
\caption{类内差异（Intra-class Variation）}
\end{subfigure}
\caption{图像分类的挑战（三）：形变和类内差异\protect\footnotemark}
\end{figure}
\footnotetext{Slides 第13--14页。}

\begin{warningbox}{像素级变化 $\neq$ 语义级变化}
即使只是将相机平移一个像素，$800 \times 600 \times 3$ 张量中的\textbf{每一个}数值都会改变。但从人类感知角度来看，图像内容完全没有变化。这正是基于原始像素的简单距离度量不适合图像分类的根本原因。
\end{warningbox}

\subsection{本章小结}

图像分类是计算机视觉中最基础、最重要的任务之一。它面临的核心挑战来自语义鸿沟——同一物体在不同视角、光照、遮挡等条件下的像素表示差异巨大。传统的基于规则（if-then-else）的方法无法扩展到复杂的现实场景，因此我们需要\textbf{数据驱动}的方法。

%% ========== Part 2: 数据驱动方法 ==========

